\subsection{Localized near-rest charged packets}

A particle rest edge and a localized detector signal concern different
preparations. They can be joined on the native kinetic and pairing source
by an explicit smooth compact packet, without introducing a spatial trap.
Let $L\ge4$, let the periodic side be $N\ge L+5$, and define
\[
 g_j=\sqrt{\frac{8}{3(L+1)}}\sin^2\frac{\pi j}{L+1}
 \quad(1\le j\le L),\qquad g_j=0\quad\hbox{otherwise}.
\]
The product $f(x)=g_{x_1}g_{x_2}g_{x_3}|\beta+,+z\rangle$
is normalized and supported in a cube $C$. With the same native
$h=K+m\beta$, $\{K,\beta\}=0$ and $E=\sqrt{m^2+K^2}$, source
fermion addition on the filled reference succeeds with vector $P_+f$,
where $P_+=(1+h/E)/2$. The charged reservoir, idle opposite reference
and full failure branch are retained.

Put $q=\pi/(L+1)$ and $D=(T-T^*)/(2i)$. Direct shifts, including both
boundary collars, give
\begin{align*}
 d_2=\|Dg\|^2&=\frac{\sin^2(2q)}3+
                    \frac{4\sin^4q}{3(L+1)},\\
 d_4=\|D^2g\|^2&=\frac8{3(L+1)}
 \left[\frac L8\sin^4(2q)-\frac12\sin^2(2q)\cos(2q)\sin^2q
                                  +\frac14\sin^4q\right].
\end{align*}
Product structure gives $a_2=\langle K^2\rangle_f=3d_2$ and
$a_4=\langle K^4\rangle_f=3d_4+6d_2^2$. Their exact finite formulas
supply rational upper bounds by enclosing $\pi$ and the elementary
trigonometric inequalities. No large spectral diagonalization is required.

The success weight is $w=(1+m\langle E^{-1}\rangle_f)/2$ and
$1-w\le\eta=a_2/(4m^2)$. For $\eta<1$, the conditional packet
$\psi=P_+f/\sqrt w$ obeys
\[
 0\le\langle E\rangle_\psi-m\le d=
   \frac{a_2}{2m(1-\eta)},\qquad
 \langle(E-m)^2\rangle_\psi\le b^2=
   \frac{a_4}{4m^2(1-\eta)}.
\]
Although $P_+$ gives the compact seed a spatial tail, its complete
outside-cube probability satisfies
\[
 \|(1-Q_C)e^{-iht}\psi\|^2\le
 \left(\sqrt{\frac\eta{1-\eta}}+|t|b\right)^2.
\]
With $L=2^{30}$, $N=4(L+1)$ and the light native charged coefficient,
the conditional mean energy is at most $3.014\,10^{-9}$ relatively above
its rest edge. The electron anchor therefore gives a localized mean energy
immediately above $0.00051099895069\,\mathrm{GeV}$, within about
$0.003014$ parts per million. Through source time $10^9$, the conditional
outside probability is below $1.103\,10^{-7}$. A preparation-vector error
$\epsilon_0$ and a complete-source Hermitian perturbation of norm $\delta$
add at most $\epsilon_0+|t|\delta$ to the amplitude bound. A per-site
perturbation bound cannot replace this total norm condition.

The exact local moments are unchanged when the volume doubles at fixed
support. Complete smaller boxes execute all native position, energy and
failure records. The large member is an analytic finite-carrier certificate,
with every valley and massless spectator retained. Its full-volume evolution
is not numerically enumerated. The complete compiler includes all ancillary
outputs and arbitrary references. Exact decomposition into commuting
four-CAR factors and a prepared-state energy estimate, derived below,
control both the complete record law and the physical-energy marginal.
With $0\le H_{\rm phys}-E_G\le282N^3$ and target mean at most
$m_{\rm upper}+d$, the required cycle count is at most
$6.564\,10^{27}$ for the stated volume and time. The energy error is
at most $10^{-6}d$ and the unhalved trace error at most $10^{-12}$.
A separate lower kinetic-moment bound exceeds that error, so the compiled
mean also lies above the rest edge. This is the physical marginal energy,
not the complete enlarged compiler Hamiltonian. An efficient local
preparation, interacting electron or stationary soliton is not inferred.

\subsection{The same packet's autonomous phase and drive response}

A single coherent preparation entangles the packet and its compensating
reference with a pointer, relative to the source ground and reference
vacuum. The idle reference Hamiltonian is zero by the declared source
contract. The success and failure operators form a complete contraction
instrument on the full carrier. Let
$C(t)=\langle\psi,e^{-iht}\psi\rangle$. Two coherent return quadratures,
their orthogonal complements and preparation failure have seven
original-input probabilities
\[
 p_{a,\pm}=\frac w8(1+|C|^2\pm2c_a),\quad
 p_{a,\perp}=\frac w4(1-|C|^2),\quad p_{\rm fail}=1-w,
 \qquad c_X=\Re C,\quad c_Y=-\Im C.
\]
Since the same spectral law has $E-m\ge0$, it gives
\[
 |C(t)-e^{-imt}|\le |t|d,\qquad |C(t)|\ge1-|t|d,
 \qquad
 \left|-\partial_t\arg C(t)-m\right|\le\frac d{1-|t|d}
\]
when $|t|d<1$. The displayed packet preserves visibility above $0.9997728$
for about twelve thousand rest cycles. Its phase and energy share one
Hamiltonian, preparation, ground and reference.

A nonzero reference cost $\mu$ multiplies the return by $e^{-i\mu t}$.
Dephasing the initial arms removes their quadrature contrast. A differential
native scalar drive $v(t)Q_{\rm source}$ multiplies it by
$\exp(i\int v)$, whereas a common source/reference gauge change cancels.
The actual full-generator inverse returns the complete state. These
controls establish prepared autonomous phase coherence of the declared
quadratic target Hamiltonian. Passive onset is not established, and the
differential scalar drive shifts its frequency. The finite phase-cycle
compiler uses continuing controls; its sampled complete-record bound
does not establish an uncontrolled microscopic oscillation or
instantaneous frequency rigidity between pulses.

\subsection{Native-current conversion into transverse radiation}
A declared finite quantum extension couples the actual charged Dirac current
$-\tau_3\otimes\sigma_r$ to four transverse cochain radiation modes.
All twelve source CAR coordinates and complete finite oscillator occupations
are retained. Total charge and lattice momentum are conserved. For a scalar
rest particle--hole pair on the filled free reference, both crossed-polarization
two-photon matrix elements of $(H/m)^2$ have magnitude
$2\sqrt{2/5}\,\eta^2$, with opposite phases. The source derivation sums
all intermediate modes and both emission orders; no annihilation projector
is inserted. A full-sector norm bound gives positive original radiation
probability above $4.3451063\,10^{-38}$ at dimensionless time $10^{-8}$.

At time fifty, complete sparse evolutions give approximately $0.03544$
and $0.02865$ for caps one and two, retaining a material cutoff difference.
Vacuum, zero-coupling and frequency-detuning controls and the interaction
energy are retained. Number-readout work is distinguished from conserved
premeasurement energy. The supplied box resonance, coupling and free
preparation define this finite experiment; physical particle identity,
continuum annihilation rates and asymptotic outgoing radiation are not
inferred from the conversion probability.

\subsection{Angular momentum and the localized field torque}

A rotational regulator separates spatial angular momentum from an
assigned spin label. Let $P_n$ retain every three-dimensional oscillator
state of total excitation at most $n$, with $n\ge2$. It commutes with
rotations, and its coordinate wavefunctions are localized Gaussian
polynomials. Define
\[
 p_i=P_n(-i\partial_i)P_n,\qquad
 L_i=P_n(-i\epsilon_{ijk}x_j\partial_k)P_n,\qquad
 J_i=L_i+\Sigma_i/2.
\]
The differential $L_i$ obey the exact rotation algebra and rotate $p_i$ as
a vector. Products of separately truncated coordinates and momenta have a
different boundary term. The native sixteen-row Clifford representation
retains four complete Dirac blocks, both signed branches and every spatial
shell. A continuous nontrivial spatial rotation cannot act by permutations
on a finite cubic site set, so exact cubic covariance alone does not
supply this continuous generator.

Retain the three real transverse field amplitudes $b$ in
$A_b(x)=e^{-|x|^2/(2R^2)}(b\times x)/2$ and their conjugate momenta $\pi$.
Gaussian integration of the actual Maxwell energy gives
\[
 c_0=\int |A_{\hat b}|^2=\frac{\pi^{3/2}R^5}{4},\qquad
 c_1=\int |\nabla\times A_{\hat b}|^2=
       \frac{5\pi^{3/2}R^3}{8},\qquad
 H_f=\frac{\kappa|\pi|^2}{2c_0}+\frac{c_1|b|^2}{2\kappa}.
\]
For the declared positive kernel
$K(b)=|\beta m+\alpha\cdot(p-qA_b)|+\Omega\mathcal N$,
the oscillator energy trap is distinct from the cubic scalar mass well.
Its actual current is $-\partial_b\langle K\rangle$, evaluated by the
complete signed spectral derivative. Rotational covariance and the
Maxwell mode equations give
\[
 \frac{d\langle J\rangle}{dt}
       =b\times\partial_b\langle K\rangle,
 \qquad
 \frac{d(b\times\pi)}{dt}
       =-b\times\partial_b\langle K\rangle.
\]
Thus $\langle J\rangle+b\times\pi$ and the total energy are conserved.
The finite executable checks these identities, complete angular records
and two integration tolerances. The tighter trajectory has angular drift
below $4.9\,10^{-10}$ and energy drift below $3.9\,10^{-9}$; these are
numerical diagnostics separate from the conservation theorem. This is a
declared rotational source with semiclassical transverse fields, without
longitudinal closure or a transfer of the cubic or interacting ground.

\subsection{A common static and precession magnetic limit}

Use the same complete-shell energy at $n=2$, $\Omega=1$, $q=-1$,
$R=6/5$ and $\kappa=m^{-6}$, with quasistatic minimization of the field.
The field treatment differs from evolving its canonical momenta, while
the matter kernel, current and energy are identical. Write
$D=\beta m+A$. The native anticommutation relation gives
$|D|=\sqrt{m^2+A^2}$ on both signs. At zero probe its full forty-row
kinetic square $T=A^2$ obeys
\[
 T(T-5I/2)=0.
\]
For the uniform magnetic probe derivative $E_B=\partial_B A$,
$Q=E_B^2$ obeys $Q(Q-I/4)(Q-I/2)=0$. These exact finite identities give
$\|A\|<8/5$, $\|E_B\|<4/5$ and $\|\partial_B T\|<3$.
On the oscillator vacuum the two native spin states have
$\langle\partial_B T\rangle=\pm1$. The magnetic coefficient follows
from the squared coupled Dirac operator, without an inserted Pauli term.

The number gap isolates the lowest axial eigenline. A positive Sylvester
inverse bounds the magnetic derivative's difference from
$(\partial_B T)/(2m)$, while the field-current Hessian makes the axial
self-consistency map a contraction. The stationary branch has exact
$J_z=1/2$ and field momentum zero. At zero self-field its representation
also has $J^2=3/4$. The fixed axial self-field can mix different $j$
values, so the dressed state's complete $J^2$ probabilities are retained
without asserting that same Casimir. The separate canonical-spin
expectation is not substituted for the angular generator.

For $B=1/10$ prepare the equal superposition of the two actual lowest
axial eigenlines and evolve through $T_{\rm obs}=64m$. All other modes
remain in the source. A full complementary Schur bound controls their
energy split; a uniform Duhamel estimate controls the state-dependent
self-field and the observed angular phase. Static and precession
normalizations use the same empty-vacuum-relative stationary energy,
including its field cost. At $m=16384$ the rational bounds give
\[
 |g_{\rm static}-2|<0.001832,\qquad |g_{\rm precession}-2|<0.004863.
\]
The bounds tend to zero along the declared mass/coupling route.
Numerical values near $1.99999999884$ and $2.00000000061$ are evaluations,
not narrow spectral or physical uncertainty intervals. This finite
low-momentum theorem supplies no physical electron mass, fine-structure
constant, anomalous moment or cubic-source continuum limit. Its
quasistatic dynamics is semiclassical, without a linear quantum-channel
claim for arbitrary references.

\subsection{Autonomous native-current attraction with flux circulations}

An explicit specialization supplies a complete autonomous charge-pair
experiment without position pinning. Retain the full native periodic
cube carrier of side four, but enable only the electron-singlet nearest
axial terms on the line $(0,0,z)$. Their literal kinetic coefficients are
three $i\sqrt2/8$ entries and the periodic-boundary $-i\sqrt2/8$.
Other hopping and Yukawa coefficients are zero by declaration. Add
positive matter number energy $2N_{\rm matter}$, scalar excitation energy,
reference energy $4N_{\rm ref}$ and the three-dimensional Green energy
$(\kappa/2)\rho_Y^T G\rho_Y$ of the actual native current. The unit
nonblank-link penalty is absent.

The complete 119-state electron-singlet matrix coordinate and its adjoint
preserve three states with outgoing hypercharges $0,-1,+1$. The selected
reference fills all thirty modes except one electron-singlet mode, giving
hypercharge $-1$ and number energy $116$. One active particle has charge
$+1$. All 30 matter modes per cell, scalar factors, reference occupations
and link complements remain in the defined source. Enumerating the native
Gauss equation
\[
 \delta_{z,y}-\delta_{z,0}+f_z-f_{z-1}=0
\]
on the four-link ring gives nine position/circulation states for each
spin. The full native column action closes this eighteen-dimensional
physical sector exactly; its complement is retained. The interacting
unit-Yukawa Hamiltonian does not preserve it, so no interacting ground
or clock is transferred.

With $J_z$ the actual link-phase derivative, the autonomous clock obeys
\[
 i[H,N_z]=J_{z-1}-J_z,\qquad
 \left(\frac{dJ_z}{dt}\right)_{\rm electric}=i[V,J_z].
\]
Here $z$ labels a bond and $J_z$ is its charge current, distinct from the
angular generator above. The Green energies at unit coupling are
$0,21/128,31/160,21/128$ around the ring. A specified coherent input on
positions one and two gives the exact inward electric contribution
$-19\sqrt2/5120$ to the outward current acceleration. Complete
position, spin, circulation and energy records follow from the same
Hamiltonian. At time eight the origin probability is about $0.47953$,
compared with $0.41818$ in the zero-electric-coupling control. Larger
coupling need not increase arrival monotonically because interference
and the finite flux cutoff are retained. Removing circulation states
changes the time-twenty position law by total variation about $0.25986$.
This is a finite quantum current response, without a classical trajectory,
physical coupling, radiation field or far-field continuum inference.

\subsection{Exact neutral momentum flatness and local preparation}
\label{sec:exact-neutral-flatness}

For the periodic side-two fermion/scalar/link source, let $Q_i$ be
fermion parity on the coordinate-one plane and let $R_i$ exchange and
invert the paired axis-$i$ links. For single-link inversion $\mathcal I$,
the full native coordinate identity
$\mathcal I U_{ab}\mathcal I=U_{ba}^\dagger$ and endpoint-current exchange
give $[Q_iR_i,H]=0$. Translation satisfies
\[
 T_j(Q_iR_i)T_j^*=(-1)^{N_f\delta_{ij}}Q_iR_i.
\]
In every odd-fermion sector these symmetries exchange translation
characters. Every energy is therefore unitarily identical at all eight
momenta, to all orders in the hopping. This includes the isolated neutral
band. Its dressed location projectors commute exactly with $H$. A local
band preparation has zero probability at other atomic locations for
every time, with at most $(v/d)^2$ in its retained virtual-profile
complement. This is an exact periodic side-two obstruction to motion,
without a transfer to larger or antiperiodic members or a moving-particle
mass interpretation.

The same neutral band has a useful preparation probability when its
input is built from known local eigenspaces. A uniform choice of neutral
cell and seven ground spins, with eight explicit local unitary services,
blank links and ten retained classical bits, prepares the reduced atomic
density $P/1024$. The local lines are algebraic projectors of actual
sixteen-dimensional physical blocks. Supplied reservoirs and local
services are declared inputs, rather than a free reset from an arbitrary
state. The recipe archive and parity-balancing reference are retained.

A finite degree-10182 polynomial $M(H)$ has $\|M-E\|\le10^{-16}$ for
the actual interacting band projector $E$. With
$\eta=\|E-P\|\le v/d$, put
\[
 A=PM^2P,\qquad q_0=(1-\eta^2)(1-10^{-16})^2,\qquad
 L=\sqrt{q_0}\,MP\frac{3P-A}{2}.
\]
The positive inverse-square-root series proves
$L^*L\le q_0P\le I$ and success probability above
$0.99999999999639$, with failure square root retained. The conditional
reduced source density is within $2.000001\,10^{-16}$ in unhalved trace
norm of $E/1024$, giving an all-time stationarity defect below
$4.000002\,10^{-16}$. It does not assert stationarity of individual
recipes or source/archive correlations. Three polynomial evaluations
require 30546 formal source-Hamiltonian applications; their finite
Stinespring services are declared, without a primitive gate compiler.
All eight momentum outcomes have probability one eighth. The complete
factorized preparation, energy, position and local-profile instrument
has 410624 records, including every failure and profile complement.
This preparation does not remove the exact flatness of the same source.

\subsection{Volume and sampling controls for the colour coupling}

The same native-seam-derived pure-colour Wilson action at $\beta=6$
is sampled on periodic four-dimensional sides $6,8,10$. Each volume
retains independent hot and cold histories of four hundred records, with
its declared warm-up, stride and final configuration. The two larger-box
chains add eight hundred records; all twenty-four hundred original
records enter the common analysis. The five rectangle means and finite
Gaussian normalization define the same three coupling readouts at lattice
lengths $1/2,3/2,5/2$.

Common-block nonlinear jackknives retain cross-scale covariance. Block
lengths $5,10,20,40,80$, split histories and hot/cold differences expose
sampling dependence. The exact derivative of each logarithmic loop ratio
gives its influence history; an initial-positive monotone covariance-pair
sum supplies an autocorrelation diagnostic. Independent-box contrast
covariances are sums of the two estimates. These finite histories do not
certify mixing or rigorous confidence coverage. The receipt retains all
volume contrasts, including effects visible at the shortest loop.

At side ten the three estimates are $1.564128$, $2.759170$ and
$4.163746$, with twenty-record-block standard errors $0.000337$,
$0.003190$ and $0.026478$. The side-ten minus side-eight contrasts are
$0.002268$, $0.023749$ and $0.027022$, with independent-box standard
errors $0.000602$, $0.005469$ and $0.046618$. The first two contrasts
remain visible relative to these sampling diagnostics; agreement of the
longest-loop estimate does not certify the absence of volume bias. All
three measured coupling values increase with the loop scale on side ten.

The comparison changes volume at fixed bare action and spacing. It tests
the finite-scale coupling trend without inserting perturbative running,
while supplying no continuum beta function, matched step-scaling limit,
particle-core radius or physical GeV conversion.
